\begin{table}[H]
\caption{Métodos computacionales identificados en los estudios que responden a RQ2.\label{tab:rq2-metodos}}
%\isPreprints{\centering}{%
\begin{adjustwidth}{-\extralength}{0cm}
%}
%\isPreprints{\begin{tabularx}{\textwidth}{CCC}}{%
\begin{tabularx}{\fulllength}{p{0.30\fulllength} X p{0.08\fulllength}}
%}
\toprule
\textbf{Método computacional} & \textbf{Estudios designados} & \textbf{Total} \\
\midrule

RAG y recuperación contextual/semántica & \cite{Khadem2025,Hadjichristofi2026,Ersoy2025,Awadid2025,Xu2025,Gijre2024,Guan2024,Su2025,Omrani2024,Mukherjee2025,Qiao2024,Chuang2024,Lakshmanan2025,Negoita2025,Thanh2025} & 15 \\

Grafos de conocimiento, ontologías y razonamiento semántico & \cite{Paneque2024,Hosseini2025,Gao2021,VergaraVargas2024,Zhang2025b} & 5 \\

Recuperación guiada por relevancia sobre bases de conocimiento & \cite{Jin2025} & 1 \\

Razonamiento basado en casos (CBR) & \cite{Jung2023} & 1 \\

Optimización evolutiva/memética guiada por LLM/RAG & \cite{Salem2026} & 1 \\

Descubrimiento causal integrado con RAG/LLM & \cite{Mascia2025} & 1 \\

Composición dirigida por modelos y refinamiento basado en reglas & \cite{Gstur2024} & 1 \\

\midrule
\textbf{Total} & & \textbf{25} \\
\bottomrule
\end{tabularx}
%\isPreprints{}{%
\end{adjustwidth}
%}
\end{table}